\clase{1}{4 de Marzo}{}

\section{Preliminares}

\begin{definition}[Espacio Vectorial]
	Sean $X$ un conjunto no vacío, $\mathbb{K} = \R$ o $\C$. $X$ es un $\mathbb{K}$ espacio vectorial (ev) si está equipado con dos operaciones:
	\begin{align*}
		X \times X &\to X \\
		(x,y) &\mapsto x + y
	.\end{align*}
	\begin{align*}
		\mathbb{K} \times X &\to X \\
		(\lambda,x) &\mapsto \lambda x
	.\end{align*}
	satisfaciendo las condiciones siguientes:
	\begin{enumerate}
		\item por un lado,

		\subitem $\forall x,y \in X,\ x + y = y + x$;

		\subitem $\forall x,y,z \in X,\ (x + y) + z = x + (y + z)$;

		\subitem $\exists 0 \in X$ tal que $\forall x \in X,\ 0 + x = x$;

		\subitem $\forall x \in X,\ \exists (-x) \in X$ tal que $x + (-x) = 0$;

		\item $\forall x \in X,\ \forall \lambda,\mu \in \mathbb{K},\ (\lambda \mu) x = \lambda (\mu x)$

		\item $\forall x,y \in X,\ \forall \lambda,\mu \in \mathbb{K},\ \lambda (x + y) = \lambda x + \lambda y$, $(\lambda + \mu) x = \lambda x + \mu x$.
	\end{enumerate}
\end{definition}

\begin{eg}
	$(-1) x = -x$, $0 \cdot x = 0$.
\end{eg}

Los elementos de $X$ se llaman vectores y los de $\mathbb{K}$ se llaman escalares.

\begin{notation}
	Sean $X$ un ev, $A \subset X,\ a \in X,\ \lambda \in \mathbb{K}$. Luego,
	\begin{align*}
		a + A &= \{ a + x;\ x \in A \} \\
		\lambda A &= \{ \lambda x;\ x \in A \} \\
		a - A &= \{ a - x;\ x \in A \} \\
		A + B &= \{ x + y;\ x \in A,\ y \in B \}
	.\end{align*}
\end{notation}

\begin{note}
	$2A \neq A + A$ en general.
\end{note}

\begin{definition}
	Sean $X$ ev, $Y \subset X$. 
	\begin{enumerate}
		\item Diremos que $Y$ un subespacio (sev) vectorial de $X$ si:
		\begin{enumerate}
			\item $0_{x} \in Y$;

			\item $\forall \alpha, \beta \in \mathbb{K},\ \alpha Y + \beta Y \subset Y$.
		\end{enumerate}

		\item Diremos que $Y$ es convexo si:
		\begin{enumerate}
			\item $\forall t \in [0,1],\ tY + (1 - t) Y \subset Y$;

			\item $\forall t \in [0,1],\ \forall y_{1},y_{2} \in Y,\ ty_{1} + (1 - t) y_{2} \in Y$.
		\end{enumerate}

		\item Diremos que $Y$ es equilibrado (balanced) si $\forall \alpha \in \mathbb{K},\ | \alpha | \leq 1,\ \alpha Y \subset Y$.

		\item Diremos que $Y$ es absorbente si $\forall x \in X,\ \exists \alpha_{x} \in \mathbb{K}$ tal que $x \in \alpha A$.
	\end{enumerate}
\end{definition}

\begin{observe}
	Sea $Z \subset X$, $X$ ev.
	\begin{itemize}
		\item $\spn (Z)$ denota el ev más chivo que contiene a $Z$. Se puede describir como la familia de todas las combinaciones lineales (cl) finitas de elementos de $Z$. $\spn (Z)$ se llama ev generado por $Z$.

		\item $\hull (Z)$ denota el conjunto convexo más chico que contiene a $Z$. $\hull(Z)$ se llama envoltura convexa de $Z$.
	\end{itemize}
\end{observe}

\begin{eg}
	Sea $u \in X$.
	\begin{align*}
		\spn(u) = \begin{cases}
			\{ 0 \} \quad \text{si } u = 0 \\
			\mathbb{K} u \quad \text{recta dirigida por } u \text{ si } u \neq 0
		\end{cases}
	.\end{align*}
\end{eg}

\begin{eg}
	\begin{enumerate}
		\item $\mathbb{K}^{n}$ donde $n \in \N$.

		\item $\mcal{F}(E; \mathbb{K})$ denota el conjunto de las funciones $f \colon E \to \mathbb{K}$ donde $E \neq \varnothing$ y
		\begin{align*}
			(f + g)(x) &= f(x) + g(x) \quad \forall x \in E \\
			(\lambda f)(x) &= \lambda f(x)
		.\end{align*}

		\item Si $(E,d)$ es un espacio métrico $\mcal{C}^{0}(E; \mathbb{K})$ es un ev.

		\item Si $(E, \xi, \mu)$ es un espacio de medida $\mscr{M}(E; \mathbb{K})$ es un ev.
	\end{enumerate}
\end{eg}

\begin{definition}[espacio vectorial topológico]
	Sea $X$ un ev y $\tau$ una topología sobre $X$ tal que
	\begin{enumerate}
		\item tanto
		\begin{align*}
			X \times X &\to X \\
			(x,y) &\mapsto x + y
		,\end{align*}
		como
		\begin{align*}
			\mathbb{K} \times X &\to X \\
			(\lambda,x) &\mapsto \lambda x
		\end{align*}
		sean continuas.

		\item $\forall x \in X,\ \{ x \}$ es cerrado.
	\end{enumerate}
	En ese caso, $X$ equipado con $\tau$ se llama espacio vectorial topológico (evt).
\end{definition}
